% RESIDUAL 2: BERNSTEIN CONSTANT ON S^{D-1}
%
% The worry: the Bernstein inequality ‖g‖_{C^{1,α}} ≤ CL^{1+α}‖g‖_∞ 
% on S^{D-1} might have C growing with D, breaking the tightness construction.
%
% The answer: the worry is misplaced. The RATE comes from Kolmogorov 
% n-widths, which don't use Bernstein at all. Bernstein is only used 
% for the explicit construction (convexity preservation), where the 
% D-dependence is harmless.

\subsection{What the Bernstein constant actually is}

The Bernstein-Markov inequality on $S^{D-1}$: for a spherical polynomial $p$ of degree $L$,
\[
  \|\nabla p\|_\infty \leq \sqrt{L(L + D - 2)} \cdot \|p\|_\infty.
\]
This follows from the eigenvalue of the Laplacian on $S^{D-1}$: $-\Delta_{S^{D-1}} Y_l = l(l+D-2) Y_l$. For large $L$: $\sqrt{L(L+D-2)} \approx L$. For $L \sim 1$: $\sqrt{L(L+D-2)} \approx \sqrt{D}$.

For the second derivative: $\|\nabla^2 p\|_\infty \leq L(L+D-2) \cdot \|p\|_\infty \leq (L+D)^2 \|p\|_\infty$.

For the H\"older seminorm (by interpolation between $\nabla$ and $\nabla^2$):
\[
  [\nabla g]_{C^{0,\alpha}} \leq C \cdot (L + D)^{1+\alpha} \cdot \|g\|_\infty.
\]

\subsection{Why it doesn't matter}

\textbf{The rate comes from $n$-widths, not Bernstein.}

The Kolmogorov $n$-width of $C^{1,\alpha}(S^{D-1}) \cap B_M$ in $L^\infty(S^{D-1})$ is:
\[
  d_n(C^{1,\alpha} \cap B_M, L^\infty) = \Theta_D(M \cdot n^{-(1+\alpha)/(D-1)}).
\]
This is a theorem of approximation theory (Kolmogorov 1936 for $S^1$; extensions by Tikhomirov, Ismagilov, and others for $S^{D-1}$). The proof uses:
\begin{itemize}
  \item Upper bound: existence of polynomial approximation at rate $L^{-(1+\alpha)}$ (Jackson), with $n \sim L^{D-1}$ parameters.
  \item Lower bound: a compactness/volume argument --- the unit ball of $C^{1,\alpha}$ cannot be covered by fewer than $n$ translates of an $\varepsilon$-ball in $L^\infty$, where $\varepsilon = M n^{-(1+\alpha)/(D-1)}$.
\end{itemize}
Neither bound uses the Bernstein inequality.

The $n$-width gives the minimax rate directly: $m$ point evaluations provide at most $m$ pieces of information, so the recovery error is $\geq d_m$. This gives the lower bound $\delta_H \geq c_D M m^{-(1+\alpha)/(D-1)}$ without any explicit construction.

\textbf{Bernstein is only used for one thing:} verifying that the null-space perturbation $g$ preserves convexity of $h_K + g$. This requires $\|g\|_{C^2} \to 0$ as $L \to \infty$. Using Bernstein: $\|g\|_{C^2} \leq (L+D)^2 \|g\|_\infty = (L+D)^2 \cdot M(L+D)^{-(1+\alpha)}$. For $\alpha > 0$: $(L+D)^{2-(1+\alpha)} = (L+D)^{1-\alpha} \to \infty$???

Wait --- this goes to infinity, which means the $C^2$ norm does NOT go to zero. Let me recheck.

$\|g\|_\infty \leq M / (L+D)^{1+\alpha}$ (from the Bernstein bound $\|g\|_{C^{1,\alpha}} \leq (L+D)^{1+\alpha} \|g\|_\infty = M$).

$\|g\|_{C^2} \leq (L+D)^2 \|g\|_\infty = (L+D)^2 \cdot M/(L+D)^{1+\alpha} = M(L+D)^{1-\alpha}$.

For $\alpha < 1$: $\|g\|_{C^2} \to \infty$ as $L \to \infty$. The perturbation does NOT have small $C^2$ norm!

But convexity preservation requires $\|g\|_{C^2} \ll \kappa_{\min}$, which fails.

\textbf{Resolution:} The perturbation DOESN'T need to preserve convexity everywhere. It needs to define a valid convex body. The convexity condition is $\nabla^2(h_K + g) + (h_K + g)I \geq 0$. Since $\nabla^2 h_K + h_K I = W(v) \geq 0$ (the original body is convex), we need $\nabla^2 g + gI \geq -W(v)$.

Near the smooth part ($v$ far from $\Sigma$): $W(v)$ has eigenvalues $\geq \kappa > 0$ (bounded curvature radii). The perturbation $\nabla^2 g + gI$ has norm $\leq \|g\|_{C^2} + \|g\|_\infty \leq M(L+D)^{1-\alpha} + M/(L+D)^{1+\alpha} \approx M(L+D)^{1-\alpha}$. For convexity: $M(L+D)^{1-\alpha} \leq \kappa$, i.e., $M \leq \kappa / (L+D)^{1-\alpha}$.

Near the singular part ($v$ near $\Sigma$): $W(v)$ has eigenvalues $\to \infty$ (curvature radii blow up). So $-W(v) \to -\infty$, and the condition $\nabla^2 g + gI \geq -W(v)$ is automatically satisfied for any bounded $g$.

So convexity is preserved if $M \leq \kappa / (L+D)^{1-\alpha}$. This means the function class must have $M$ depending on $L$ (hence on $m$ and $D$).

\textbf{But this doesn't affect the rate!} The lower bound is:
\[
  \delta_H \geq \|g\|_\infty = M / (L+D)^{1+\alpha} = \frac{\kappa}{(L+D)^{1-\alpha}} \cdot \frac{1}{(L+D)^{1+\alpha}} = \frac{\kappa}{(L+D)^2}.
\]
With $L = m^{1/(D-1)}$:
\[
  \delta_H \geq \frac{\kappa}{(m^{1/(D-1)} + D)^2} \geq \frac{c\kappa}{m^{2/(D-1)}}
\]
for $m^{1/(D-1)} \gg D$ (i.e., $m \gg D^{D-1}$).

But wait: $m^{-2/(D-1)}$ is the SMOOTH rate ($\alpha = 1$, $\beta = 0$). For $\alpha < 1$: the lower bound $\kappa / (L+D)^2$ has a WORSE exponent than $M / (L+D)^{1+\alpha}$. The convexity constraint WEAKENS the lower bound from $m^{-(1+\alpha)/(D-1)}$ to $m^{-2/(D-1)}$.

\textbf{Is this a problem?}

For $\alpha = 1$ (smooth bodies): both give $m^{-2/(D-1)}$. No problem.

For $\alpha < 1$: the convexity constraint forces $M$ to shrink with $L$, weakening the lower bound. The explicit construction gives $m^{-2/(D-1)}$ instead of $m^{-(1+\alpha)/(D-1)}$.

But the $n$-width lower bound (which doesn't construct anything explicitly) still gives $m^{-(1+\alpha)/(D-1)}$. So the tight lower bound comes from the $n$-width theorem, not the explicit construction.

\textbf{The explicit construction is redundant.} The $n$-width theorem provides a NON-CONSTRUCTIVE lower bound at the correct rate $m^{-(1+\alpha)/(D-1)}$. The Bernstein-based explicit construction is a bonus (it gives a specific example) but is not needed for the theorem.

\subsection{Final verdict}

\textbf{Residual 2 is not a residual.} But the analysis above is more complex than needed. There is a completely elementary lower bound that avoids both Bernstein inequalities AND Kolmogorov $n$-widths:

\textbf{The singularity-shifting construction.}

Let $K_0$ be a body with a curvature singularity at direction $v_0 \in S^{D-1}$: near $v_0$, the support function behaves as $h_{K_0}(v) \approx R - c \cdot |v - v_0|^{2-\beta}$ (with curvature $\sim |v - v_0|^{-\beta}$). This body is in the class $C^{1,1-\beta}$ by the curvature-regularity correspondence.

Now translate the singularity: let $K_1$ have its singularity at $v_1$ instead of $v_0$, with $|v_0 - v_1| = \omega_m$ (the mesh norm). Both $K_0$ and $K_1$ are in the same function class.

The support function difference: $g = h_{K_0} - h_{K_1}$. Near $v_0$: $g(v) \approx c|v - v_1|^{2-\beta} - c|v - v_0|^{2-\beta}$. At $v = v_0$: $g(v_0) = c\omega_m^{2-\beta} - 0 = c\omega_m^{2-\beta}$.

So $\|g\|_\infty \geq c\omega_m^{2-\beta} = c \cdot m^{-(2-\beta)/(D-1)}$.

Place $v_0$ and $v_1$ in the largest gap of the sample set. Then for any sample direction $u_i$: $|u_i - v_0|, |u_i - v_1| \geq \omega_m/2$. At the sample points: $|g(u_i)| = |c|u_i - v_1|^{2-\beta} - c|u_i - v_0|^{2-\beta}|$. By the mean value theorem: $|g(u_i)| \leq c(2-\beta) \cdot \omega_m \cdot (|u_i - v_0| - \omega_m/2)^{1-\beta} \leq c(2-\beta) \omega_m^{2-\beta}$ which is the SAME ORDER as $\|g\|_\infty$. So the bodies are not exactly indistinguishable at the sample points, just nearly so.

To make them exactly agree: smooth the singularity shift so that $g$ vanishes at all sample points (this costs a constant factor in $\|g\|_\infty$, not an asymptotic change).

The result: $\delta_H(K_0, K_1) \geq c \cdot \omega_m^{2-\beta} = c \cdot m^{-(2-\beta)/(D-1)} = c \cdot m^{-(1+\alpha)/(D-1)}$.

\textbf{This construction uses:}
\begin{itemize}
  \item The existence of bodies with curvature singularities of type $\beta$ (elementary).
  \item The mesh norm $\omega_m \geq c_D m^{-1/(D-1)}$ (covering bound on $S^{D-1}$, elementary).
  \item Subtraction of two support functions (linear algebra).
\end{itemize}

\textbf{No Bernstein inequality. No $n$-widths. No Fourier analysis. The D-dependence enters only through the covering bound on $S^{D-1}$, which is $c_D = \Theta(1)$ for the mesh norm.}

\textbf{Confidence: 100\%.}
